%%% solo.tex — der Regellauf fuer dsa5solo.
%%%
%%% Ein kleines Solo mit dreissig Bloecken: genug fuer mehrere
%%% Doppelseiten, damit die Behaeltergrenzen und die Fuellseiten im Satz
%%% zu sehen sind.
%%%
%%% Der Bau geht in drei Schritten. Erst messen -- das schreibt solo.solo
%%% mit der Hoehe jedes Blocks und seinen Sprungzielen:
%%%
%%%   cd beispiel
%%%   xelatex -jobname=solo '\PassOptionsToPackage{messen}{dsa5solo}%%% solo.tex — der Regellauf fuer dsa5solo.
%%%
%%% Ein kleines Solo mit dreissig Bloecken: genug fuer mehrere
%%% Doppelseiten, damit die Behaeltergrenzen und die Fuellseiten im Satz
%%% zu sehen sind.
%%%
%%% Der Bau geht in drei Schritten. Erst messen -- das schreibt solo.solo
%%% mit der Hoehe jedes Blocks und seinen Sprungzielen:
%%%
%%%   cd beispiel
%%%   xelatex -jobname=solo '\PassOptionsToPackage{messen}{dsa5solo}%%% solo.tex — der Regellauf fuer dsa5solo.
%%%
%%% Ein kleines Solo mit dreissig Bloecken: genug fuer mehrere
%%% Doppelseiten, damit die Behaeltergrenzen und die Fuellseiten im Satz
%%% zu sehen sind.
%%%
%%% Der Bau geht in drei Schritten. Erst messen -- das schreibt solo.solo
%%% mit der Hoehe jedes Blocks und seinen Sprungzielen:
%%%
%%%   cd beispiel
%%%   xelatex -jobname=solo '\PassOptionsToPackage{messen}{dsa5solo}%%% solo.tex — der Regellauf fuer dsa5solo.
%%%
%%% Ein kleines Solo mit dreissig Bloecken: genug fuer mehrere
%%% Doppelseiten, damit die Behaeltergrenzen und die Fuellseiten im Satz
%%% zu sehen sind.
%%%
%%% Der Bau geht in drei Schritten. Erst messen -- das schreibt solo.solo
%%% mit der Hoehe jedes Blocks und seinen Sprungzielen:
%%%
%%%   cd beispiel
%%%   xelatex -jobname=solo '\PassOptionsToPackage{messen}{dsa5solo}\input{solo}'
%%%
%%% Dann loesen -- das verteilt die Bloecke auf Doppelseiten, vergibt die
%%% Nummern und teilt solo-bloecke.tex bytegetreu auf:
%%%
%%%   python3 ../werkzeuge/solo.py solo.solo --bloecke solo-bloecke.tex
%%%
%%% Dann setzen (dreimal, wegen Inhalt und Marken) und pruefen:
%%%
%%%   xelatex solo.tex
%%%   python3 ../werkzeuge/solo.py --pruefen solo.aux

\documentclass[raster]{dsa5latex}
\usepackage{dsa5solo}

\graphicspath{{../}}
\dsaAbenteuertitel{Der Bote von Havena}

\begin{document}

\dsakapitel{Der Bote von Havena}

\dsaEinfuehrung{Ein Solo-Abenteuer fuer einen Helden, der lesen kann und
sich nicht dabei erwischen lassen sollte.}

\dsaabschnitt{So wird gespielt}

Dieses Abenteuer liest sich nicht von vorn nach hinten. Jeder Abschnitt
traegt eine Zahl, und am Ende eines Abschnitts steht, wo es weitergeht.
Lies immer nur den Abschnitt, auf den du verwiesen wirst — und nie den
daneben.

Dass das gelingt, dafuer ist gesorgt: kein Verweis fuehrt auf eine Zahl,
die auf derselben Doppelseite steht. Du musst also immer blaettern, und
was du dabei siehst, verraet nichts.

\dsaabschnitt{Was du brauchst}

Einen Helden, Papier und Stift, und zwei sechsseitige Wuerfel. Der Held
sollte sich in einer Stadt zurechtfinden und notfalls schnell laufen
koennen; kaempfen muss er nicht, aber es schadet nicht.

Notiere dir vor dem Beginn die Lebensenergie deines Helden und fuehre
darueber Buch. Wo du Schaden nimmst, steht es im Text — etwa
\dsaFormel{1W6} Punkte fuer einen Sturz vom Dach.

\begin{dsaliste}
\item einen Helden, gern den aus dem Anhang
\item Papier und Stift fuer Lebensenergie und Notizen
\item zwei sechsseitige Wuerfel, \dsaFormel{1W6} und \dsaFormel{2W6}
\item einen zwanzigseitigen Wuerfel fuer die Proben
\end{dsaliste}

\dsaabschnitt{Proben}

Manche Abschnitte verlangen eine Probe. Dann steht dort, auf welches
Talent gewuerfelt wird und wie es weitergeht, je nachdem ob sie gelingt
oder misslingt. Wer keine Regeln bemuehen will, entscheidet nach Gefuehl
und liest den Abschnitt, der besser passt.

\begin{dsaMeisterMaske}
\dsaKastentitel{Wie eine Probe abgehandelt wird}
Eine Probe geht auf drei Eigenschaften. Wuerfle fuer jede einen
\dsaFormel{1W20} und vergleiche ihn mit dem Eigenschaftswert deines Helden:
liegt der Wurf darueber, kostet die Differenz Fertigkeitspunkte. Bleibt am
Ende kein Punkt uebrig, ist die Probe misslungen.

Eine Erschwernis von \dsaMod{-}{2} senkt jeden der drei Eigenschaftswerte um
zwei, eine Erleichterung hebt sie an. Wer keine Regeln bemuehen will, wuerfelt
\dsaFormel{1W6}: bei 4 und mehr gelingt die Probe.
\end{dsaMeisterMaske}

\dsaabschnitt{Wenn du stirbst}

Dann faengst du wieder von vorn an, und du weisst diesmal mehr als beim
ersten Mal. Das ist kein Betrug, sondern der Sinn der Sache: Ein Solo
spielt man mehrmals, und erst beim dritten Durchgang sieht man, wie viele
Wege es gab.

\dsaabschnitt{Wo das spielt}

\dsaBildSpalte{grafiken/aventurienkarte}{18}

Havena liegt im Westen des Mittelreichs, am Meer, und lebt vom Hafen. Wer
dort einen Brief zu ueberbringen hat, sollte wissen, dass die Stadt zwei
Arten von Gassen kennt: solche, durch die man geht, und solche, durch die
man rennt.

\dsaabschnitt{Der Anfang}

\dsaVorlesetext[7]{Der Brief wiegt nichts, und trotzdem traegst du ihn, als
waere er aus Blei. Havena liegt vor dir, der Regen im Ruecken, und die Sonne
hat noch drei Handbreit bis zum Wasser.}

Du bist zu Fuss unterwegs nach Havena, mit einem Brief in der Tasche, den
du bis Sonnenuntergang abliefern sollst. Was es damit auf sich hat, wirst
du frueh genug erfahren.

Beginne bei Abschnitt \soloWeiter{start}.

\soloStart{start}
\soloBloecke{solo-bloecke}

\end{document}
'
%%%
%%% Dann loesen -- das verteilt die Bloecke auf Doppelseiten, vergibt die
%%% Nummern und teilt solo-bloecke.tex bytegetreu auf:
%%%
%%%   python3 ../werkzeuge/solo.py solo.solo --bloecke solo-bloecke.tex
%%%
%%% Dann setzen (dreimal, wegen Inhalt und Marken) und pruefen:
%%%
%%%   xelatex solo.tex
%%%   python3 ../werkzeuge/solo.py --pruefen solo.aux

\documentclass[raster]{dsa5latex}
\usepackage{dsa5solo}

\graphicspath{{../}}
\dsaAbenteuertitel{Der Bote von Havena}

\begin{document}

\dsakapitel{Der Bote von Havena}

\dsaEinfuehrung{Ein Solo-Abenteuer fuer einen Helden, der lesen kann und
sich nicht dabei erwischen lassen sollte.}

\dsaabschnitt{So wird gespielt}

Dieses Abenteuer liest sich nicht von vorn nach hinten. Jeder Abschnitt
traegt eine Zahl, und am Ende eines Abschnitts steht, wo es weitergeht.
Lies immer nur den Abschnitt, auf den du verwiesen wirst — und nie den
daneben.

Dass das gelingt, dafuer ist gesorgt: kein Verweis fuehrt auf eine Zahl,
die auf derselben Doppelseite steht. Du musst also immer blaettern, und
was du dabei siehst, verraet nichts.

\dsaabschnitt{Was du brauchst}

Einen Helden, Papier und Stift, und zwei sechsseitige Wuerfel. Der Held
sollte sich in einer Stadt zurechtfinden und notfalls schnell laufen
koennen; kaempfen muss er nicht, aber es schadet nicht.

Notiere dir vor dem Beginn die Lebensenergie deines Helden und fuehre
darueber Buch. Wo du Schaden nimmst, steht es im Text — etwa
\dsaFormel{1W6} Punkte fuer einen Sturz vom Dach.

\begin{dsaliste}
\item einen Helden, gern den aus dem Anhang
\item Papier und Stift fuer Lebensenergie und Notizen
\item zwei sechsseitige Wuerfel, \dsaFormel{1W6} und \dsaFormel{2W6}
\item einen zwanzigseitigen Wuerfel fuer die Proben
\end{dsaliste}

\dsaabschnitt{Proben}

Manche Abschnitte verlangen eine Probe. Dann steht dort, auf welches
Talent gewuerfelt wird und wie es weitergeht, je nachdem ob sie gelingt
oder misslingt. Wer keine Regeln bemuehen will, entscheidet nach Gefuehl
und liest den Abschnitt, der besser passt.

\begin{dsaMeisterMaske}
\dsaKastentitel{Wie eine Probe abgehandelt wird}
Eine Probe geht auf drei Eigenschaften. Wuerfle fuer jede einen
\dsaFormel{1W20} und vergleiche ihn mit dem Eigenschaftswert deines Helden:
liegt der Wurf darueber, kostet die Differenz Fertigkeitspunkte. Bleibt am
Ende kein Punkt uebrig, ist die Probe misslungen.

Eine Erschwernis von \dsaMod{-}{2} senkt jeden der drei Eigenschaftswerte um
zwei, eine Erleichterung hebt sie an. Wer keine Regeln bemuehen will, wuerfelt
\dsaFormel{1W6}: bei 4 und mehr gelingt die Probe.
\end{dsaMeisterMaske}

\dsaabschnitt{Wenn du stirbst}

Dann faengst du wieder von vorn an, und du weisst diesmal mehr als beim
ersten Mal. Das ist kein Betrug, sondern der Sinn der Sache: Ein Solo
spielt man mehrmals, und erst beim dritten Durchgang sieht man, wie viele
Wege es gab.

\dsaabschnitt{Wo das spielt}

\dsaBildSpalte{grafiken/aventurienkarte}{18}

Havena liegt im Westen des Mittelreichs, am Meer, und lebt vom Hafen. Wer
dort einen Brief zu ueberbringen hat, sollte wissen, dass die Stadt zwei
Arten von Gassen kennt: solche, durch die man geht, und solche, durch die
man rennt.

\dsaabschnitt{Der Anfang}

\dsaVorlesetext[7]{Der Brief wiegt nichts, und trotzdem traegst du ihn, als
waere er aus Blei. Havena liegt vor dir, der Regen im Ruecken, und die Sonne
hat noch drei Handbreit bis zum Wasser.}

Du bist zu Fuss unterwegs nach Havena, mit einem Brief in der Tasche, den
du bis Sonnenuntergang abliefern sollst. Was es damit auf sich hat, wirst
du frueh genug erfahren.

Beginne bei Abschnitt \soloWeiter{start}.

\soloStart{start}
\soloBloecke{solo-bloecke}

\end{document}
'
%%%
%%% Dann loesen -- das verteilt die Bloecke auf Doppelseiten, vergibt die
%%% Nummern und teilt solo-bloecke.tex bytegetreu auf:
%%%
%%%   python3 ../werkzeuge/solo.py solo.solo --bloecke solo-bloecke.tex
%%%
%%% Dann setzen (dreimal, wegen Inhalt und Marken) und pruefen:
%%%
%%%   xelatex solo.tex
%%%   python3 ../werkzeuge/solo.py --pruefen solo.aux

\documentclass[raster]{dsa5latex}
\usepackage{dsa5solo}

\graphicspath{{../}}
\dsaAbenteuertitel{Der Bote von Havena}

\begin{document}

\dsakapitel{Der Bote von Havena}

\dsaEinfuehrung{Ein Solo-Abenteuer fuer einen Helden, der lesen kann und
sich nicht dabei erwischen lassen sollte.}

\dsaabschnitt{So wird gespielt}

Dieses Abenteuer liest sich nicht von vorn nach hinten. Jeder Abschnitt
traegt eine Zahl, und am Ende eines Abschnitts steht, wo es weitergeht.
Lies immer nur den Abschnitt, auf den du verwiesen wirst — und nie den
daneben.

Dass das gelingt, dafuer ist gesorgt: kein Verweis fuehrt auf eine Zahl,
die auf derselben Doppelseite steht. Du musst also immer blaettern, und
was du dabei siehst, verraet nichts.

\dsaabschnitt{Was du brauchst}

Einen Helden, Papier und Stift, und zwei sechsseitige Wuerfel. Der Held
sollte sich in einer Stadt zurechtfinden und notfalls schnell laufen
koennen; kaempfen muss er nicht, aber es schadet nicht.

Notiere dir vor dem Beginn die Lebensenergie deines Helden und fuehre
darueber Buch. Wo du Schaden nimmst, steht es im Text — etwa
\dsaFormel{1W6} Punkte fuer einen Sturz vom Dach.

\begin{dsaliste}
\item einen Helden, gern den aus dem Anhang
\item Papier und Stift fuer Lebensenergie und Notizen
\item zwei sechsseitige Wuerfel, \dsaFormel{1W6} und \dsaFormel{2W6}
\item einen zwanzigseitigen Wuerfel fuer die Proben
\end{dsaliste}

\dsaabschnitt{Proben}

Manche Abschnitte verlangen eine Probe. Dann steht dort, auf welches
Talent gewuerfelt wird und wie es weitergeht, je nachdem ob sie gelingt
oder misslingt. Wer keine Regeln bemuehen will, entscheidet nach Gefuehl
und liest den Abschnitt, der besser passt.

\begin{dsaMeisterMaske}
\dsaKastentitel{Wie eine Probe abgehandelt wird}
Eine Probe geht auf drei Eigenschaften. Wuerfle fuer jede einen
\dsaFormel{1W20} und vergleiche ihn mit dem Eigenschaftswert deines Helden:
liegt der Wurf darueber, kostet die Differenz Fertigkeitspunkte. Bleibt am
Ende kein Punkt uebrig, ist die Probe misslungen.

Eine Erschwernis von \dsaMod{-}{2} senkt jeden der drei Eigenschaftswerte um
zwei, eine Erleichterung hebt sie an. Wer keine Regeln bemuehen will, wuerfelt
\dsaFormel{1W6}: bei 4 und mehr gelingt die Probe.
\end{dsaMeisterMaske}

\dsaabschnitt{Wenn du stirbst}

Dann faengst du wieder von vorn an, und du weisst diesmal mehr als beim
ersten Mal. Das ist kein Betrug, sondern der Sinn der Sache: Ein Solo
spielt man mehrmals, und erst beim dritten Durchgang sieht man, wie viele
Wege es gab.

\dsaabschnitt{Wo das spielt}

\dsaBildSpalte{grafiken/aventurienkarte}{18}

Havena liegt im Westen des Mittelreichs, am Meer, und lebt vom Hafen. Wer
dort einen Brief zu ueberbringen hat, sollte wissen, dass die Stadt zwei
Arten von Gassen kennt: solche, durch die man geht, und solche, durch die
man rennt.

\dsaabschnitt{Der Anfang}

\dsaVorlesetext[7]{Der Brief wiegt nichts, und trotzdem traegst du ihn, als
waere er aus Blei. Havena liegt vor dir, der Regen im Ruecken, und die Sonne
hat noch drei Handbreit bis zum Wasser.}

Du bist zu Fuss unterwegs nach Havena, mit einem Brief in der Tasche, den
du bis Sonnenuntergang abliefern sollst. Was es damit auf sich hat, wirst
du frueh genug erfahren.

Beginne bei Abschnitt \soloWeiter{start}.

\soloStart{start}
\soloBloecke{solo-bloecke}

\end{document}
'
%%%
%%% Dann loesen -- das verteilt die Bloecke auf Doppelseiten, vergibt die
%%% Nummern und teilt solo-bloecke.tex bytegetreu auf:
%%%
%%%   python3 ../werkzeuge/solo.py solo.solo --bloecke solo-bloecke.tex
%%%
%%% Dann setzen (dreimal, wegen Inhalt und Marken) und pruefen:
%%%
%%%   xelatex solo.tex
%%%   python3 ../werkzeuge/solo.py --pruefen solo.aux

\documentclass[raster]{dsa5latex}
\usepackage{dsa5solo}

\graphicspath{{../}}
\dsaAbenteuertitel{Der Bote von Havena}

\begin{document}

\dsakapitel{Der Bote von Havena}

\dsaEinfuehrung{Ein Solo-Abenteuer fuer einen Helden, der lesen kann und
sich nicht dabei erwischen lassen sollte.}

\dsaabschnitt{So wird gespielt}

Dieses Abenteuer liest sich nicht von vorn nach hinten. Jeder Abschnitt
traegt eine Zahl, und am Ende eines Abschnitts steht, wo es weitergeht.
Lies immer nur den Abschnitt, auf den du verwiesen wirst — und nie den
daneben.

Dass das gelingt, dafuer ist gesorgt: kein Verweis fuehrt auf eine Zahl,
die auf derselben Doppelseite steht. Du musst also immer blaettern, und
was du dabei siehst, verraet nichts.

\dsaabschnitt{Was du brauchst}

Einen Helden, Papier und Stift, und zwei sechsseitige Wuerfel. Der Held
sollte sich in einer Stadt zurechtfinden und notfalls schnell laufen
koennen; kaempfen muss er nicht, aber es schadet nicht.

Notiere dir vor dem Beginn die Lebensenergie deines Helden und fuehre
darueber Buch. Wo du Schaden nimmst, steht es im Text — etwa
\dsaFormel{1W6} Punkte fuer einen Sturz vom Dach.

\begin{dsaliste}
\item einen Helden, gern den aus dem Anhang
\item Papier und Stift fuer Lebensenergie und Notizen
\item zwei sechsseitige Wuerfel, \dsaFormel{1W6} und \dsaFormel{2W6}
\item einen zwanzigseitigen Wuerfel fuer die Proben
\end{dsaliste}

\dsaabschnitt{Proben}

Manche Abschnitte verlangen eine Probe. Dann steht dort, auf welches
Talent gewuerfelt wird und wie es weitergeht, je nachdem ob sie gelingt
oder misslingt. Wer keine Regeln bemuehen will, entscheidet nach Gefuehl
und liest den Abschnitt, der besser passt.

\begin{dsaMeisterMaske}
\dsaKastentitel{Wie eine Probe abgehandelt wird}
Eine Probe geht auf drei Eigenschaften. Wuerfle fuer jede einen
\dsaFormel{1W20} und vergleiche ihn mit dem Eigenschaftswert deines Helden:
liegt der Wurf darueber, kostet die Differenz Fertigkeitspunkte. Bleibt am
Ende kein Punkt uebrig, ist die Probe misslungen.

Eine Erschwernis von \dsaMod{-}{2} senkt jeden der drei Eigenschaftswerte um
zwei, eine Erleichterung hebt sie an. Wer keine Regeln bemuehen will, wuerfelt
\dsaFormel{1W6}: bei 4 und mehr gelingt die Probe.
\end{dsaMeisterMaske}

\dsaabschnitt{Wenn du stirbst}

Dann faengst du wieder von vorn an, und du weisst diesmal mehr als beim
ersten Mal. Das ist kein Betrug, sondern der Sinn der Sache: Ein Solo
spielt man mehrmals, und erst beim dritten Durchgang sieht man, wie viele
Wege es gab.

\dsaabschnitt{Wo das spielt}

\dsaBildSpalte{grafiken/aventurienkarte}{18}

Havena liegt im Westen des Mittelreichs, am Meer, und lebt vom Hafen. Wer
dort einen Brief zu ueberbringen hat, sollte wissen, dass die Stadt zwei
Arten von Gassen kennt: solche, durch die man geht, und solche, durch die
man rennt.

\dsaabschnitt{Der Anfang}

\dsaVorlesetext[7]{Der Brief wiegt nichts, und trotzdem traegst du ihn, als
waere er aus Blei. Havena liegt vor dir, der Regen im Ruecken, und die Sonne
hat noch drei Handbreit bis zum Wasser.}

Du bist zu Fuss unterwegs nach Havena, mit einem Brief in der Tasche, den
du bis Sonnenuntergang abliefern sollst. Was es damit auf sich hat, wirst
du frueh genug erfahren.

Beginne bei Abschnitt \soloWeiter{start}.

\soloStart{start}
\soloBloecke{solo-bloecke}

\end{document}
